% !TeX root = main.tex
\section{Conclusions}
\label{sec:Conclusions}

The evolution of scientific software is a requirement to continue fulfilling sustainability aspects of science.
We have shown how the strategic expansion of certain modern software engineering practices can impact the overall development process of the QMCPACK framework targeting HPC systems.
We provide empirical data as evidence for how expanding CI efforts with a combination of cloud-based CI on free GitHub Actions runners and project self-maintained hardware resources for testing on NVIDIA and AMD GPUs has enabled QMCPACK readiness as exascale systems continue to be deployed.
In addition, the targeted quality strategy adds extra checks and refactoring efforts to sanitize critical code paths, track code testing coverage, and allow for multiple supported build configurations.
We outlined some of the important software engineering efforts in the QMCPACK software evolution roadmap: the removal of legacy code implementations, adding more unit tests to track coverage, modernizing current checkpoint-restart and I/O capabilities, refactoring compile-time configuration into runtime variants.
Our goal is to showcase how the targeted, predictive, proactive (rather than reactive), and selective application of modern software engineering practices have benefited a well-established HPC scientific code and community like QMCPACK.
We continue to advocate in the research software engineering (RSE) community that understanding the ``critical path" for a scientific software project is the first step in the successful and pragmatic integration of software engineering and scientific computing.
Publishing these types of experiences enriches the RSE community as a whole by showcasing the value and impact of their work.
